\chapter{细胞呼吸：电子怎样驱动 ATP 合成}

\begin{kaichang}
现在，请你站起来，原地快跑一分钟，或者做三十个深蹲。做完之后停下来，注意自己的身体：你在大口喘气，脸在发烫，可能已经开始出汗，心脏咚咚直跳。

好，坐下，我们趁这个状态出一道题。你刚才吸进去的那一大批氧气，最后都变成了什么？

先猜。绝大多数人的第一反应是：“变成了二氧化碳，又被我呼出去了。”毕竟体育课上就是这么讲的：吸进氧气，呼出二氧化碳。把你的猜测写在纸上。这一章结束时我们对答案——正确答案会让很多人意外：你吸进的氧，绝大多数变成了\textbf{水}。而二氧化碳里的氧，另有来历。

从喘气到出汗这一整套身体反应，背后是一台藏在你几乎每个细胞里的发电站。这一章，我们就拆开这台发电站，看电子怎样像下楼梯一样逐级释放能量，又怎样驱动 ATP 的合成。
\end{kaichang}

\section{不冒火苗的燃烧}

先看几个谁都能观察到的现象。

第一，你不能停止呼吸。人可以几天不吃饭、一天不喝水，但完全断氧几分钟，大脑就开始不可逆地受损。说明身体里有某种过程对氧气的需求是\textbf{持续}的，不能“存起来慢慢用”。

第二，运动越剧烈，喘气越快。安静坐着时你每分钟呼吸十几次，冲刺之后能到三四十次。身体像一个会自动调档的发动机：活儿越重，吸气越多。

第三，运动时身体发热。冬天里跑两圈就暖和了，这不是错觉——你的体温调节系统在拼命散热：皮肤血管扩张（所以脸红）、出汗。

第四，呼出的气是潮湿的。对着冷玻璃窗呼一口气，玻璃上立刻蒙一层水雾。有人会说：水嘛，当然是喝进去又排出来的。但请看一个数字：一个成年人即使一整天不喝水，呼气带走的水汽也不会停止。

第五，一个精确的测量：把 1 克葡萄糖放在弹式量热计里烧掉，放出的热量，和一只动物把 1 克葡萄糖彻底“用掉”放出的总热量，是一样的——约 \ud{16}{kJ}。殊途同归，起点和终点的能量差分毫不差。可烧葡萄糖有火苗，你身体里没有任何火苗。同一笔能量，火焰是一次性全倒出来，身体却把它分成了几十次小额支付，还把大约三分之一存进了 ATP 这张“储蓄卡”（第 50 章）。

这中间的支付系统是怎么工作的？镜头得推进到细胞里面。

\begin{shiyan}
\textbf{材料}：一面小镜子或一块冷玻璃（冰箱里的冰镇饮料瓶也行）、秒表或手机计时器。

\textbf{步骤}：（1）安静坐五分钟，数自己一分钟的呼吸次数，记下。（2）原地快跑一分钟，立刻再数一分钟呼吸次数，记下。（3）对着镜子近距离缓慢呼气，观察镜面。

\textbf{观察点}：两次呼吸次数差了多少倍？镜面上的水雾从哪里来？注意：呼气的水雾并不全是“喝进去又路过”的水——等你读完本章，你会知道其中一部分水是在你细胞里刚刚“出厂”的。

\textbf{安全提示}：运动量力而行，头晕就停下；本章其他内容涉及线粒体抑制剂等毒物，\textbf{严禁}在家尝试任何相关操作，只能看老师演示或正规视频。
\end{shiyan}

\section{丙酮酸进了发电站}

回忆第 52 章：葡萄糖在细胞质里经过糖酵解，被拆成两个丙酮酸，净赚 2 个 ATP，顺手把几个高能电子存进了电子载体 NADH。这笔买卖其实很不划算——葡萄糖的大部分能量还锁在丙酮酸里，就像拆开包裹只拿走了赠品，主件还没动。

丙酮酸的主件要去哪儿拆？去\textbf{线粒体}。先认识一下这个细胞器的结构，因为本章所有机关都装在它的膜上：线粒体有\textbf{两层膜}。外膜平缓，像厂房的外墙；内膜反复向内折叠，形成一排排叫“嵴”的褶皱，像暖气片——折叠是为了在有限体积里塞进尽可能大的膜面积。内膜围着的内腔叫\textbf{基质}，内膜与外膜之间的窄缝叫\textbf{膜间隙}。请记住这三个地名：基质、内膜、膜间隙，接下来它们要反复出现。

顺便说，线粒体不是每个细胞平均分一个。哪里的活儿重，哪里的发电站就多：你的心肌细胞一刻不停跳一辈子，体积的三分之一被线粒体占据；精子的尾巴里密密排着一圈线粒体，专供游泳的鞭打；而成熟红细胞干脆一个线粒体也没有——它把全部空间让给血红蛋白，自己只靠糖酵解那 2 个 ATP 凑合过日子。你身体里的发电站布局，就是一张精确的“能耗地图”。

丙酮酸穿过两层膜进入基质后，先做一次“进站安检”：被切下一个碳，以 \chem{CO_2} 的形式放掉——注意，这就是你呼出的二氧化碳的出生地之一。剩下的二碳碎片挂到一个搬运工（辅酶 A）身上，随后进入一个环形流水线：\textbf{柠檬酸循环}（也叫三羧酸循环或克雷布斯循环）。

这条流水线的设计很妙：它不是一条有头有尾的生产线，而是一个\textbf{圈}。二碳碎片进站，和一个四碳分子拼成六碳的柠檬酸；机器转一圈，沿途把两个碳以 \chem{CO_2} 踢出去，到终点时恰好把起点的四碳分子原样造回来，迎接下一批进站原料。所以柠檬酸既是第一个产物，又循环再生——这也是它名字的由来。

一圈转下来，真正的收获不是 ATP（直接到手的只有 1 个 ATP 等价物），而是\textbf{电子}：流水线沿途有四道工序，把高能电子从碳骨架上剥下来，装进电子载体——3 个 NADH 和 1 个 \chem{FADH_2}（它和 NADH 是同类角色，第 50 章介绍过这批“电子运钞车”）。加上进站安检那一步又产生的 1 个 NADH，一个丙酮酸彻底拆掉，共交出 4 个 NADH、1 个 \chem{FADH_2} 和 3 个 \chem{CO_2}。

一个葡萄糖出两个丙酮酸，数字翻倍。到此为止，葡萄糖的碳骨架已经\textbf{全部}变成 \chem{CO_2} 离场，可细胞到手的 ATP 只有寥寥 4 个。能量去哪儿了？绝大部分变成了 NADH 和 \chem{FADH_2} 身上的高能电子。碳走了，电子留下了。真正的发电，现在才开始。

\section{电子下楼梯，质子上山坡}

高能电子的兑现场所，是镶嵌在线粒体\textbf{内膜}上的一串巨型蛋白质机器：\textbf{电子传递链}。它主要由四个蛋白复合体组成，编号 I 到 IV，像一排逐级下降的台阶。

NADH 把电子交到第一级台阶（复合体 I），电子随后依次传给 III、IV。为什么要排成台阶，而不是一步到底？这正是第 27 章能量规则的生物版应用：电子从高能级落向低能级，每落一级只释放一小份自由能，蛋白质机器才有机会把这一小份\textbf{接住、用掉}；要是一步跳到底，能量一次性以热的形式散掉，就像把水桶从十楼直接扔下去——动能全变成了巨响。分步走，是细胞版的“水逐级推动一串小水车”。

那各级台阶把接住的能量用来干什么？答案出人意料：用来\textbf{搬质子}。复合体 I、III、IV 每传递一轮电子，就利用释放的能量把质子（\chem{H^+}）从基质泵过内膜，送到膜间隙。电子在下楼，质子却在上坡——电子下楼释放的能量，被转存成了质子的“高度差”。

结果是膜两侧积累出巨大的不平衡：膜间隙里质子又多、又带正电，基质里质子稀少。第 50 章说过，浓度差和电荷差都是储存能量的方式，两股力量合起来叫\textbf{质子动力势}——它就像水库大坝两侧的水位差，只差 \ud{0.2}{V} 左右，可膜只有几纳米厚，折算成电场强度高达每厘米几十万伏，比高压输电线周围的电场还猛。

大坝要有水轮机才能发电。线粒体的水轮机是第五台蛋白机器：\textbf{ATP 合酶}，也嵌在内膜上。它是目前已知最小的旋转马达：膜间隙的质子从它的通道里涌回基质，推动一个转子高速旋转（每秒上百圈），转子的扭动让另一端的催化部位变形，把 ADP 和磷酸“压”成 ATP。第 50 章提醒过，ATP 的能量不在某根“高能键”里，而在于水解产物远更稳定、浓度远离平衡——这里我们看到的是它的另一面：ATP 是\textbf{被质子的洪流顶着造出来的}，只要梯度在，这台马达就不停转，一个 ATP 合酶每秒能造上百个 ATP。

整条链还有最后一块拼图：那些“下了楼”的电子，最后去哪？总得有谁接收它们，否则链条很快就会堵死。接收者站在台阶最末端——就是你吸进来的\textbf{氧气}。氧在复合体 IV 处接过电子，再从基质里抓两个质子，结合成水：
\[\chem{O_2} + 4\chem{H^+} + 4e^- \rightarrow 2\chem{H_2O}\]
氧的角色是\textbf{最终电子受体}，一个永远敞开的电子出口。它不当燃料、不直接碰葡萄糖，只负责在终点把用完的电子接走。这就是为什么断氧几分钟就致命：出口一关，电子全线堵死，质子泵停摆，梯度消散，ATP 合酶停转——全身的发电站同时拉闸。

\section{把整套流程写成符号}

到这一节，符号才登场。先给整个过程一个总账本。一个葡萄糖彻底氧化，总反应式和你熟悉的燃烧式一模一样：
\[\chem{C_6H_{12}O_6} + 6\chem{O_2} \rightarrow 6\chem{CO_2} + 6\chem{H_2O}\]
但请注意：这个式子只写了起点和终点，中间几十步全部省略了。燃烧是一步跳到底，细胞呼吸是把同一笔能量分成几十次支付。式子相同，路径天差地别——方向由热力学决定，路径和速度由酶决定，这是第一册反复强调的老规则在生命里的兑现。

电子的“运钞车”交接，可以写成：
\[\chem{NADH} \rightarrow \chem{NAD^+} + \chem{H^+} + 2e^-\]
卸完货的 \chem{NAD^+} 空车开回糖酵解和柠檬酸循环，继续接活——这就是为什么第 52 章说细胞必须不断再生 \chem{NAD^+}，否则连糖酵解都会熄火。

整条电子传递链的能量景观，可以画成这样一张“楼梯图”：

\begin{center}
\begin{tikzpicture}
\draw (0,3.4) rectangle (2,4.2); \node at (1,3.8) {\small 复合体 I};
\draw (2.6,2.5) rectangle (4.6,3.3); \node at (3.6,2.9) {\small 复合体 III};
\draw (5.2,1.6) rectangle (7.2,2.4); \node at (6.2,2.0) {\small 复合体 IV};
\draw (7.8,0.6) rectangle (9.8,1.4); \node at (8.8,1.0) {\small \chem{O_2}$\rightarrow$\chem{H_2O}};
\draw[->,thick] (-0.2,4.6) -- (0.6,4.2); \node at (0.4,4.8) {\small NADH 进站};
\draw[->,thick] (2,3.3) -- (2.6,3.0);
\draw[->,thick] (4.6,2.4) -- (5.2,2.1);
\draw[->,thick] (7.2,1.5) -- (7.8,1.1);
\draw[->,thick] (1,3.4) -- (1,2.6) node[below]{\small 泵 \chem{H^+}};
\draw[->,thick] (3.6,2.5) -- (3.6,1.7) node[below]{\small 泵 \chem{H^+}};
\draw[->,thick] (6.2,1.6) -- (6.2,0.8) node[below]{\small 泵 \chem{H^+}};
\end{tikzpicture}
\end{center}

图旁必须说明：这是人为的表示法。复合体的真实形状是缠成一团的蛋白质，大小、间距、颜色都不代表真实外观；它只表达一件事——电子逐级下坡，每一级带动一台质子泵。

\begin{qiaoliang}
来算总账。一个葡萄糖的全部“电子收成”：糖酵解出 2 个 NADH，丙酮酸进站出 2 个，柠檬酸循环出 6 个 NADH 加 2 个 \chem{FADH_2}。电子传递链的“汇率”大约是：每个 NADH 的电子最终换来约 2.5 个 ATP，每个 \chem{FADH_2} 约 1.5 个（它从第二级台阶进站，少泵一轮质子，所以汇率低）。于是：$10\times 2.5 + 2\times 1.5 = 28$ 个 ATP 来自电子传递链，加上糖酵解和柠檬酸循环直接给的 4 个，总计约 30～32 个 ATP。对比一下：如果氧气缺席，细胞只能靠糖酵解拿到 2 个。有氧呼吸把能量提取率提高了约 15 倍——这就是大型多细胞生物（比如你）付得起大脑账单的化学前提。

再估一个数量级，感受 ATP 的周转有多疯狂。一个成年人即使整天躺着，每天也要合成并消耗约 40～70 千克 ATP——接近你的体重。但你全身任何时刻只存着几十克 ATP。结论：ATP 不是仓库里的存货，而是流水线上每秒都在印、每秒都在花的现钞。你安静坐着，体内的 ATP 合酶大军每秒也在合成数以百万亿计的 ATP 分子。
\end{qiaoliang}

\section{一个被嘲笑了十几年的想法}

\begin{zhengju}
今天我们说“质子梯度驱动 ATP 合成”，语气平淡得像常识。可在 1961 年，这个想法被主流学界当成怪谈。

当时的主流理论叫\textbf{化学偶联假说}：既然糖酵解是用“高能化学物质”一步步传递能量，那么电子传递链多半也是先造出某种神秘的高能中间体（人们叫它 X\raisebox{0.5ex}{\texttildelow}P 之类），再由它去造 ATP。这是个合理的类比推理，而且给出了明确的任务：把这个中间体分离出来。于是全世界许多实验室找了十几年——什么也没找到。一个理论预言的核心物质始终不露面，这本身就该让人警惕。

1961 年，英国生化学家彼得·米切尔提出了一个完全不同的方案，叫\textbf{化学渗透学说}：根本没有什么高能中间体。电子传递链真正的产品是把质子泵过膜，\textbf{膜两侧的质子梯度本身就是储能形式}；质子经由 ATP 合酶回流时，梯度能直接变成 ATP。这个假说要求膜承担一个全新的角色——不再是装东西的袋子，而是大坝和水库。

反应几乎是敌意。当时的教科书说能量只能存在化学键里，米切尔的理论被讥为“薄膜玄学”；他在会议上宣讲时，同行当面驳斥。更麻烦的是他的处境：身体不好、经费紧张，后来干脆搬到乡间，把一座庄园改建成私人实验室（格林研究所），靠养牛场和自家积蓄维持研究。

米切尔没有靠雄辩翻盘，他设计了一批\textbf{可被判死刑的实验}。关键的一个逻辑是：如果化学偶联是对的，那么破坏膜的完整性最多影响效率；如果化学渗透是对的，那么膜一旦漏质子，ATP 合成就该彻底停摆——而电子传递本身照常进行。实验结果清清楚楚：把膜弄漏，电子传递照常、ATP 停产。化学偶联假说无法自然地解释这个结果。

更漂亮的证据在 1974 年到来。斯托克尼乌斯和拉克尔做了一个“拼装”实验：把一种细菌的光驱动质子泵（细菌视紫红质）和 ATP 合酶装进\textbf{同一个}人造脂质小泡，小泡里不放任何细胞提取物。一照光，泵开始搬质子，小泡里随即出现 ATP。这个系统里除了泵、膜、梯度和 ATP 合酶，什么都没有——所谓“高能化学中间体”根本没有藏身之处。再配上一个老悬案的破解：某些物质（如 2,4-二硝基酚）能让呼吸照常却产热不产 ATP，化学渗透学说一句话解释——它们是把质子偷偷运过膜的“走私犯”，梯度被漏光了，能量全部变成热。

1978 年，米切尔独自获得诺贝尔化学奖。这个故事最值得记住的，不是“坚持就会胜利”——科学史上坚持错误想法的人更多——而是\textbf{翻盘的筹码}：一个理论即使在少数派地位，只要它比对手更准确地预言了“漏膜实验”和“重组小泡”的结果，它就该赢。证据到位之前，被嘲笑十几年是科学共同体的常态，不是阴谋。
\end{zhengju}

\section{这个模型什么时候不够用}

化学渗透图景极其强大，但它也有边界，说清楚边界才算真懂。

第一，线粒体不是必需品。细菌没有线粒体，把电子传递链和 ATP 合酶装在\textbf{细胞膜}上，照样用质子梯度发电。这说明化学渗透是比细胞器更基本的原理——事实上，线粒体自己就起源于远古细菌，这是后面章节要讲的故事。

第二，质子梯度不只造 ATP。许多细菌直接用质子回流驱动鞭毛旋转，像用同一座水库同时供电和供水；一些物质跨膜运输也蹭这股梯度。“梯度 = ATP 专用电池”只是主要用途，不是唯一用途。

第三，电子传递和 ATP 合成不是焊死的。褐色脂肪细胞的线粒体内膜上有一种解偶联蛋白，故意给质子开一条旁路，让梯度直接变成热——新生儿和冬眠动物靠它保暖。所以“呼吸的全部意义是产 ATP”这个说法，在你的肩胛骨之间就有一片反例。

第四，氧不是唯一的电子出口。许多微生物用硝酸盐、硫酸盐甚至铁离子当最终电子受体（这叫无氧呼吸）。它们活得好好的，只是电子落差更小、能量收成更低——氧之所以是王牌受体，正因为它“接电子”的欲望（电负性，第 11 章）几乎最强，给电子留出的下坡最长。

第五，30～32 个 ATP 是个约数。膜总有一点漏电，搬运 NADH 进出线粒体要花手续费，真实产量因细胞类型和状态浮动。生物学里的定量数字几乎都带着这种“约”字。

\begin{wuqu}
“你吸进的氧气，最后变成你呼出的二氧化碳。”——流传极广，而且错得很有教育意义。拆开看：\chem{CO_2} 是在丙酮酸进站和柠檬酸循环中从\textbf{食物的碳骨架}上剥下来的，它的氧原子来自葡萄糖和水，跟你吸进的那口 \chem{O_2} 没有直接关系；你吸进的 \chem{O_2} 在电子传递链末端接住电子和质子，变成的是\textbf{水}。最直接的证据来自同位素标记实验（第 6 章讲过这种追踪术）：让人吸入用氧-18 标记的氧气，标记出现在身体产生的水里，而不是呼出的 \chem{CO_2} 里。这也顺便解释了那个总反应式为什么“骗人”：式子两边的氧原子根本不是同一批。原子的旅程，要追踪了才知道。
\end{wuqu}

\section{回到那三十个深蹲}

现在可以重新解释开场那一分钟身体发生的一切了。你的肌肉细胞在深蹲时疯狂消耗 ATP，线粒体开足马力：柠檬酸循环加速剥电子，电子传递链全速运电子，质子洪流冲进 ATP 合酶。电子出口的需求量暴涨，所以呼吸加深加快——你喘的气，是在给亿万条电子传递链\textbf{补充终点接收站}，同时把柠檬酸循环放出的 \chem{CO_2} 运走。

发热呢？电子传递链的耦合本就不是百分之百，每传一轮电子都有一部分自由能直接以热散失；更何况肌肉收缩本身也只有约四分之一的能量变成机械功。几十万亿个线粒体同时开工，余热让你的体温调节系统拉响警报：血管扩张（脸红）、出汗。你对着镜子呼出的水雾里，就混着电子传递链末端刚刚造出来的水——你吸进的氧，此刻正在以水的形式离开你的身体。

对答案时间：氧气变成了水，二氧化碳来自食物骨架。猜对了吗？

\section{从毒药到减肥药：这条链的脆弱与强悍}

理解了一台机器，就理解了它会怎样坏掉。氰化物（\chem{CN^-}）之所以是剧毒，是因为氰根离子恰好能卡进复合体 IV 的血红素铁中心，把氧气挡在门外。电子传不到终点，整条链停摆——诡异的是，中毒者血液里氧气充足，细胞却在“富氧窒息”，皮肤反而呈现异常的潮红，因为血里的氧根本用不掉。一氧化碳是另一类打法：它抢先霸占血红蛋白（第 48 章），让氧气运不到组织。两种毒，一个堵终点站，一个断运输线，都指向同一个命门：电子必须有出口。

反过来的教训更惊心。20 世纪 30 年代，有人发现 2,4-二硝基酚这种“质子走私犯”能让呼吸的能量全变成热，于是把它当减肥药卖——脂肪确实被快速烧掉，可剂量稍过，体温就失控飙升，多名使用者死于高热。它是药理学史上最黑暗的章节之一：机制理解正确，剂量窗口却窄到杀人。这也是本书反复说的红线——\textbf{剂量、证据等级、个体差异}，在“代谢类神药”的广告面前，请永远先想这条电子传递链。

医学也用这条链做好事：常用降糖药二甲双胍的作用之一，就是轻度抑制复合体 I，让细胞能量状态发生微妙变化，从而调节全身代谢；许多杀虫剂和抗生素的靶点，也是某种电子传递链的变体。而发热这个“副产品”，在演化眼里根本不是废物：恒温动物把线粒体的余热当成了维持体温的主力热源，你之所以能在冬天不穿羽绒服也活下来，部分要感谢电子传递链的“效率不高”。

最后埋一个伏笔。这一章的故事里，电子从食物流向氧，放出能量；能不能反过来——用外界的能量，把电子从水里硬拽出来、推上坡，安到二氧化碳身上，逆着造出糖？这需要一台能把光变成电子流的机器，以及把撕开水时释放的氧气处理掉的本事。植物恰好两样都有，而且它们发电站的核心——质子梯度加 ATP 合酶——和线粒体几乎一模一样，连证据故事都遥相呼应。下一章，第 54 章，我们去叶绿体里看这台反向运转的机器，顺便回答：你呼出的水，和树叶放出的氧，是怎样在同一个化学循环里遥遥相望的。

\begin{tiaozhan}
本框可跳过，不影响主线。质子动力势有两笔账：电势差 $\Delta\psi$ 和浓度差 $\Delta\mathrm{pH}$，合起来写作 $\Delta p = \Delta\psi - 59\,\mathrm{mV}\cdot\Delta\mathrm{pH}$，线粒体里总值约 \ud{220}{mV}。把这份能量摊到电子上：一个 NADH 送出 2 个电子，三级泵共搬约 10 个质子过膜；而 ATP 合酶每造 1 个 ATP 大约放 3 个质子回流，再加上把 ADP 和磷酸运进基质要花约 1 个质子的梯度——所以每个 NADH 约值 $10/4 = 2.5$ 个 ATP，这就是正文“汇率”的来历。\chem{FADH_2} 从复合体 II 进站，绕过了第一级泵，少搬 4 个质子，所以只值约 1.5 个。再验证一次效率：1 摩尔葡萄糖燃烧放能约 \ud{2870}{kJ}，30 摩尔 ATP 水解能约 $30\times 30.5 \approx 915$ \ud{}{kJ}，存下约三分之一——剩下的三分之二哪去了？你冬天不觉得冷，就是答案。
\end{tiaozhan}

\section*{练一练}

\begin{sikaoti}
\item 画一张“物质流 + 能量流”双线图：跟踪一个葡萄糖分子的碳（走到 \chem{CO_2}）、电子（走到 \chem{O_2}）和能量（走到 ATP 和热）三条路线，标出每条路线经过的“站点”（细胞质、线粒体基质、内膜、膜间隙）。
\item 氰化物中毒者的血液明明富含氧气，细胞却会“窒息”，且皮肤呈异常潮红。用电子传递链的术语解释这三个现象各自的直接原因。
\item 假如用一种化学物质在线粒体内膜上戳出许多允许质子自由通过的小孔，预测：（a）电子传递速率会上升、下降还是不变？（b）ATP 产量如何变化？（c）体温如何变化？这套推理能解释 2,4-二硝基酚减肥药的致命事故吗？
\item 给一名志愿者吸入用氧-18 标记的 \chem{^{18}O_2}，预测标记元素\textbf{最先}大量出现在呼出的 \chem{CO_2} 里，还是体液与呼出水汽的水里？说明理由，并指出这个实验推翻了哪个流行说法。
\item 细菌没有线粒体，却也能进行有氧呼吸。它们把电子传递链和 ATP 合酶装在哪里？这个事实和“线粒体起源于细菌”的假说之间有什么关系？
\item 估算：一个成年人每天静息合成约 \ud{50}{kg} ATP。ATP 的摩尔质量约 \ud{507}{g\cdot mol^{-1}}，换算一下你的细胞每天大约合成多少摩尔、多少个 ATP 分子（用 \ud{6.02\times 10^{23}}{mol^{-1}}）。再想想：既然产量如此惊人，为什么“补 ATP 口服液”对健康人基本是智商税？
\end{sikaoti}
